% Authoritative LaTeX manuscript; Markdown is a generated compatibility projection.

% <!-- part-title: Alignment -->

% <!-- chapter-title: Where Obligations Can Be Enforced -->

\chapter{Where Obligations Can Be Enforced}

The preceding section distinguished what governance mechanisms do. Their placement determines what work they can examine and which actions they can still affect. A permission may act on one attempted operation; a completion check may need the combined evidence from several changes. We now examine these boundaries and how to choose among them.

\section{Intervention Boundaries}\label{heading:reasoning-input-action-work-unit-and-artifactruntime}

% <!-- point: enforcement-can-act-at-several-boundaries-choose-where-the-obligation-is-visible | An obligation can be enforced at several boundaries; choose where it is visible and a mechanism can enforce it. | terms: intervention-boundary -->

Where the environment intervenes determines what it can examine and control. \autoref{intervention-boundaries} distinguishes reasoning input, where guidance ordinarily operates, from three enforcement boundaries: individual actions, composed units of work, and the resulting artifact or runtime state. Moving toward larger units of work can make additional obligations legible because more context and evidence are available.

\begin{figure}[htbp]
\centering
\includegraphics[width=\linewidth]{book/latex/assets/intervention-boundaries-1ac2c40a9b.pdf}
\caption{\emph{Intervention boundaries.} Guidance influences reasoning; enforcement constrains actions or determines whether work and artifacts are accepted. Moving from individual actions to larger units of work can make additional obligations legible as evidence accumulates. Later boundaries are not necessarily stronger. Choose the earliest boundary where the obligation can be evaluated and enforced.}
\label{intervention-boundaries}
\end{figure}

% <!-- section-terms: intervention-boundary, reasoning-input, action-boundary, work-unit-boundary, artifact-runtime-boundary -->

% <!-- point: the-first-boundary-is-reasoning-input | The first location is reasoning input: ordinarily guidance, not enforcement. | terms: reasoning-input, intervention-boundary -->

% <!-- index-def: reasoning-input -->

The first location is \textbf{reasoning input}: the context, representations, examples, and instructions put in front of the reasoner. This is where Modeling and context engineering do much of their work: a finite reasoner benefits when relevant system knowledge is externalized and selected for the task. But placement in context is ordinarily guidance, not enforcement. The agent can misread, forget, or disregard what it was shown. Guidance helps an agent avoid mistakes. Enforcement can prevent those mistakes from being accepted.



% <!-- point: the-second-boundary-is-the-action-one-attempted-side-effect | The first enforcement boundary is the action: one attempted side effect the harness can inspect or refuse. | terms: action-boundary, intervention-boundary -->

% <!-- index-def: action-boundary -->

The first of the three \emph{enforcement} boundaries is the \textbf{action boundary} exposed by the harness and its tools. Here, one attempted side effect is in view: a tool call, an API request, a shell command, or a file mutation.

% <!-- box-family: inset -->

% <!-- inset-domain: AGENT HARNESS -->

% mage-box-prefix: Inset —

\begin{bookbox}{inset}{Hooks and Permissions}
\label{inset-hooks-and-permissions}
Many agent harnesses let engineers intervene in an agent's work through permissions and lifecycle hooks. For example, a deployment tool can be permitted to inspect production state while its credentials deny permission to modify it. A pre-action hook can invoke a check on a proposed command; if the check is configured to block violations, the harness refuses the command. A post-action hook can check the result or collect evidence, but the operation has already occurred. That evidence can inform a later acceptance decision: a CI or deployment gate can reject a larger unit of work when required checks fail.

We can understand these facilities by separating \emph{where} an intervention occurs from \emph{what} the configured mechanism does. The boundary determines what evidence is available in time to affect the decision. The mechanism determines how that evidence is used: the same failed check at the same hook might produce a warning or stop execution. If the deployment tool must not modify production state, permissions enforced by the production service can refuse a write before it takes effect. A warning after a write cannot satisfy that obligation: the forbidden effect has already occurred. The engineer therefore starts with what must be prevented or established, then chooses a boundary with sufficient evidence and a mechanism whose response enforces that requirement.
\end{bookbox}

% <!-- point: the-third-boundary-is-the-work-unit-a-meaningful-unit-of-many-actions | The second enforcement boundary is the work unit, where whole-work obligations become legible. | terms: work-unit-boundary, intervention-boundary -->

% <!-- index-def: work-unit-boundary -->

The second enforcement boundary is the \textbf{work-unit boundary}, where many actions compose into a meaningful unit: a completed task, pull request, build, migration step, or deployment candidate. Obligations that no single action can express become legible here: that the model and the implementation agree at task completion, that a change carries all its required evidence, that "done" actually holds. A pre-commit check can reject a local defect; a task-completion audit can weigh the whole body of work; a deployment gate can refuse a candidate whose accumulated evidence is inadequate. The larger scope makes obligations involving several actions available for evaluation, even when no individual action provides enough evidence.

% <!-- point: the-fourth-boundary-is-the-artifact-or-runtime-forbidden-state-made-unexpressible | The third enforcement boundary is the artifact or runtime: forbidden state made impossible to express or accept. | terms: artifact-runtime-boundary, intervention-boundary, ban-lint -->

% <!-- index-def: artifact-runtime-boundary -->

The third enforcement boundary is the \textbf{artifact or runtime boundary}: the resulting system or running state. Types, narrow APIs, architectural seams, schemas, runtime permissions, and admission policies can make a forbidden state impossible to express or impossible to accept. For example, a service account may have permission to read production records but no permission to change them. The permission restricts what the running system can do, regardless of how the agent wrote its code.

\section{The Semantic Gap and Earliest Decidable Boundary}\label{heading:the-semantic-gap}

% <!-- section-terms: semantic-gap, earliest-legible-boundary, model-drift -->

% <!-- point: an-obligation-can-be-perfectly-stated-and-unenforceable-at-the-boundary-you-check | An obligation can be stated clearly and still be unenforceable at the boundary where you try to check it. | terms: semantic-gap, earliest-legible-boundary -->

Which of these boundaries provides enough evidence? We may know exactly what to check but try to check it too soon. The required evidence may span several changes, or the property may emerge only when components operate together. A semantic gap occurs when the chosen boundary does not expose enough information to decide whether the obligation holds.

The end-to-end argument in system design explains why this information boundary matters. A function that needs end-to-end knowledge cannot in general be implemented completely at a lower layer that lacks that knowledge~\citep{saltzerreedclark1984e2e}. That argument concerns where to place functionality in a distributed system; the corresponding question here is where to place enforcement over engineering work. An enforcing mechanism needs access to the information required to decide whether the obligation holds.

\autoref{semantic-gap} illustrates two such cases. Although both arise during software development, they require checks at different boundaries.

\begin{figure}[htbp]
\centering
\includegraphics[width=\linewidth]{book/latex/assets/semantic-gap-levels-bf0468495b.pdf}
\caption{\emph{The semantic gap.} Different obligations become decidable at different boundaries. Model--implementation correspondence may be required at task completion, while rendered-page correctness requires integrated execution. Place each check at the earliest boundary where sufficient evidence is available.}
\label{semantic-gap}
\end{figure}

% <!-- point: model-drift-at-the-commit-is-the-wrong-granularity | Checking model drift at the commit is the wrong granularity; the property belongs to the completed unit of work. | terms: model-drift, semantic-gap -->

% <!-- index-def: semantic-gap -->

% <!-- index-example: model-drift -->

Consider an agent changing a system's call graph. The architectural model must agree with the implementation when the task is complete, but several intermediate commits may legitimately leave them inconsistent. Requiring correspondence after every commit would reject valid intermediate work. At task completion, the accumulated code changes and model updates can be checked together. The obligation belongs to the completed task, not to each commit.

% <!-- point: a-deployment-can-pass-every-check-and-still-serve-a-blank-page | Test the served page in a browser at the earliest boundary where the relevant composition can be exercised; a successful HTTP fetch alone does not establish correct rendering. | terms: semantic-gap, earliest-legible-boundary -->

A different obligation requires the assembled system to behave correctly. A deployment can pass source, build, and API checks yet serve a blank page. Those checks establish useful properties, but not that the application displays the expected content. A browser smoke test can load the page and check its rendered output. The test belongs at the earliest boundary where the integrated application can be exercised, such as a preview or staging deployment. A successful HTTP response alone does not establish that the page rendered correctly.

% <!-- gloss: Epic | A tracked unit of work larger than one commit: a feature or effort with its own file, plan, and definition of done, carried out over several agent dispatches. -->

% <!-- gloss: Definition of done (DoD) | The explicit list of conditions a unit of work must meet to count as finished — tests present, docs updated, the model in sync — checked at the end rather than assumed. -->

% <!-- section-terms: earliest-legible-boundary, semantic-gap -->

% <!-- point: place-enforcement-at-the-earliest-boundary-where-the-obligation-is-legible-and-decidable | The design rule: place enforcement at the earliest boundary where the obligation is both legible and decidable. | terms: earliest-legible-boundary, semantic-gap -->

% <!-- index-def: earliest-legible-boundary -->

Both examples illustrate the same rule for placing enforcement:

\begin{quote}
\textbf{Place enforcement at the earliest boundary where the obligation is both legible and decidable.}
\end{quote}

\emph{Earliest} matters because delay makes a failure more expensive to repair. \emph{Legible} matters because checking before the necessary evidence exists produces false refusals or meaningless results. Static type checks can reject code that violates declared typing rules before execution; action hooks can govern individual tool calls; task-completion checks can evaluate obligations spanning several commits; and deployment gates can evaluate accumulated evidence before release. Other properties become observable only later, such as a leak emerging across repeated calls or a cost overrun visible across many jobs. No earlier boundary can establish what the available evidence does not yet reveal.

% <!-- point: the-part-two-rhyme-coarsest-representation-earliest-boundary | The rhyme: Modeling chooses the reduction that answers the question without carrying unnecessary detail; Alignment chooses the earliest boundary where the obligation is legible. | terms: semantic-gap, earliest-legible-boundary, modeling-principle, alignment-principle -->

\magevalue{chapter:modeling} and \magevalue{chapter:alignment} therefore use companion selection rules:

% <!-- inline-quote -->

\begin{quote}
\textbf{Modeling} --- \emph{choose the reduction that makes the engineering \textbf{question} answerable without carrying unnecessary detail.}

\textbf{Alignment} --- \emph{choose the earliest boundary at which the \textbf{obligation} becomes legible and decidable.}
\end{quote}

Both rules keep irrelevant detail from driving the engineering decision.

\begin{itemize}
\tightlist
\item
  Modeling avoids carrying detail the question does not need.
\item
  Alignment avoids enforcing an obligation before the necessary evidence exists.
\end{itemize}

Choose the wrong abstraction in \magevalue{chapter:modeling} and the required property is not represented. Choose the wrong boundary here and the required evidence is not available.

% <!-- point: real-boundaries-are-sometimes-approximate-a-roughly-right-hook-beats-the-perfect-one | Real boundaries are sometimes approximate: a hook that fires early enough to leave budget for the work beats the theoretically ideal hook. | terms: earliest-legible-boundary, semantic-gap -->

Real boundaries sometimes require a margin for the mechanism to act. Suppose an agent must preserve a durable hand-off before its working context is compacted. Waiting until no room remains may leave too little capacity to write it. A practical trigger can act earlier, when context is becoming full and the durable record is stale. The trigger approximates the desired boundary while preserving enough time and capacity to complete the required action.

An enforceable obligation needs more than a model and a checker. The obligation must be independent of the work being judged, the evidence must cover the relevant behavior, and the check must operate where that evidence is available. These requirements also leave engineers a separate decision about whether the result should control admission.

\section{Placement, Strength, and Feedback}\label{heading:placement-and-strength-are-different-decisions}

% <!-- section-terms: intervention-boundary, guidance-vs-enforcement -->

% <!-- point: the-four-boundaries-do-not-form-a-monotonic-scale | The four boundaries do not form a monotonic weak-to-strong scale. | terms: intervention-boundary, guidance-vs-enforcement -->

A later check is not necessarily a stronger check. A narrowly scoped API permission may enforce its obligation more strictly than a late CI check; a type may rule out one state completely while saying nothing about a system-level obligation a task-completion gate can weigh. \textbf{Where a mechanism acts and how strongly it enforces an obligation are different design decisions.}

% <!-- point: substantial-enforcement-can-exist-without-rich-models | Substantial enforcement can exist without rich models; richer models expand the properties the environment can state and govern. | terms: guidance-vs-enforcement, modeling-principle -->

Richer models can expand the properties the environment is able to state and govern, but substantial enforcement can exist without them. A bounded execution sandbox, for example, can constrain actions without requiring a rich model of the system it contains.

% <!-- section-terms: intervention-boundary, feedback -->

% <!-- point: every-boundary-determines-where-feedback-can-return | Boundary placement affects what feedback can return and when; timely evidence lets an agent use the reasoning and context already at hand, while delayed feedback may require reconstruction. | terms: intervention-boundary, feedback -->

% <!-- index-def: intervention-boundary -->

% <!-- index-def: engineered-agent-loop -->

The boundaries also determine what feedback can return and when. An action hook can explain why a call was refused; a CI check can report a violated invariant in the change under review; a deployment validator can identify failed checks or missing required evidence. Placing a check where the obligation can first be evaluated allows useful evidence to return while the relevant work is still in progress. The agent can then revise the change using the reasoning and context already at hand. If feedback arrives after it has moved on, it may have to reconstruct what it changed and why before it can understand the failure and correct the work.

% <!-- point: these-mechanisms-do-not-act-independently-deferred-to-3-5 | These mechanisms do not act independently; the interplay is its own problem, deferred to §3.5. | terms: intervention-boundary -->

Controls within the GEE can also interact or conflict; \magevalue{sec:when-guardrails-collide} treats their growing portfolio as a system in its own right.

\section{Enforcement Boundaries in Practice}\label{heading:enforcement-boundaries-in-practice}

The appropriate boundary depends on the obligation and the evidence needed to evaluate it. DocAble uses different boundaries to govern discovery and completed work; Docker and Spotify illustrate how other environments govern outbound requests and submission for review.

% <!-- worked-examples: alignment-principle -->

\subsection{DocAble: Dynamic Context Injection and Discovery}\label{heading:docable-dynamic-context-injection-and-discovery}

% Evidence: E19 | source: book/data/local-evidence/evidentiary-basis-11202026/integration/records/E19.md | kind: dated-case-evidence
DocAble's engineering environment must supply agents with the rules relevant to their tasks without requiring them to reconstruct every applicable obligation from the repository. It must also distinguish supplying those rules from establishing that required discovery occurred. Dynamic context injection addresses the first problem; a separate gate addresses the second. The selected knowledge includes architectural invariants, ownership declarations, lint requirements, and links to design decisions.

% Evidence: E19 | source: book/data/local-evidence/evidentiary-basis-11202026/integration/records/E19.md | kind: dated-case-evidence
% Evidence: E58 | source: book/data/local-evidence/evidentiary-basis-11202026/integration/records/E58.md | kind: dated-case-evidence
DocAble selects context twice. At dispatch, an agent hook matches the work order's anticipated files against file-scoped constraint declarations and the models' file-to-component mappings. The resulting knowledge joins the standing instructions in the agent's initial context. After planning, the agent triggers a second extraction using its refined file set. This selection may surface obligations missed initially or discard material no longer relevant to the proposed change. This is \emph{dynamic context injection}: the knowledge supplied to an agent changes as the work becomes better understood. \autoref{docable-context-selection} illustrates both selections and the separate gate that requires evidence of the second.

\begin{figure}[htbp]
\centering
\includegraphics[width=\linewidth]{book/latex/assets/deliver-context-two-layer-15d7b97afa.pdf}
\caption{\emph{Dynamic context injection and enforcement.} A dispatch hook selects initial engineering knowledge from anticipated files and components. After planning, the agent selects context again using its refined scope. Standing instructions remain available throughout. A Git precommit gate requires evidence of the second selection for the initial commit that permits implementation to begin.}
\label{docable-context-selection}
\end{figure}

% Evidence: E19 | source: book/data/local-evidence/evidentiary-basis-11202026/integration/records/E19.md | kind: dated-case-evidence
% Evidence: E39 | source: book/data/local-evidence/evidentiary-basis-11202026/integration/records/E39.md | kind: dated-case-evidence
Before implementation begins, the agent must make an initial commit demonstrating that the required discovery occurred. A separate Git pre-commit gate rejects that commit if the dispatch has no recorded second-extraction event. Its condition is deliberately narrower than ``the agent followed the rules.'' It does not compare the extraction's file set with the committed files, or establish that the selected constraints were complete, correct, read, or followed. The gate requires a discovery event; correspondence checks and implementation gates must evaluate the obligations that discovery exposes.

\subsection{DocAble: Completion Audits}\label{heading:docable-completion-audits}

% Evidence: E19 | source: book/data/local-evidence/evidentiary-basis-11202026/integration/records/E19.md | kind: dated-case-evidence
% Evidence: E20 | source: book/data/local-evidence/evidentiary-basis-11202026/integration/records/E20.md | kind: mechanism-and-practice
% Evidence: E39 | source: book/data/local-evidence/evidentiary-basis-11202026/integration/records/E39.md | kind: dated-case-evidence
Completion requires a different view of the work. DocAble's definition-of-done review asks a reviewer agent to rerun the task's tests and lints against the completed work rather than trust earlier status records. A separate closing gate checks that the declared commits are present on the main branch. Several intermediate commits may legitimately leave a model and its implementation temporarily inconsistent; the completed task is where their required agreement can be assessed. The two mechanisms therefore answer different questions: did the required discovery occur, and does the finished work carry the evidence needed for acceptance? Neither question is answered merely by placing instructions in context. ({[}appendix: appendix-a-governed-delegation{]} develops a blueprint for model selection, work-order construction, and completion checks.)

\subsection{Other Environments: Docker and Spotify}\label{heading:example-docker}

Docker Sandboxes routes outbound HTTP and HTTPS requests through a proxy on the host. Under a configured deny-by-default policy, the proxy checks the destination domain against an allowlist and blocks requests to unapproved domains. For example, suppose an agent tries to upload a diagnostic log to a domain absent from the allowlist. The proxy blocks that request before it reaches the destination, even if the agent considers the upload useful. This enforces the domain restriction; it does not establish that content sent to or obtained from an allowed domain is safe~\citep{mikhalev2026untrusted}.

Spotify requires a proposed change to pass the relevant verifiers before its background coding system opens a pull request. The environment selects independent verifiers from repository contents: finding \texttt{pom.xml}, for example, activates a Maven verifier. The agent can invoke verification while editing, and a Claude Code stop hook runs all relevant verifiers before PR creation. If one fails, the system returns an error and does not open the PR. The agent can therefore obtain feedback during its work, while the final invocation makes verification a condition of submitting that work for review~\citep{charas2025feedback}.

\magevalue{chapter:software-factories} examines how such controls fit into software factories, including Spotify's migration system.

% <!-- worked-examples-end -->

We can now state an obligation, choose machinery suited to it, and place that machinery where the necessary evidence is available. But we cannot anticipate every obligation before engineering begins. New behavior, failures, and changes in requirements will expose decisions that the environment does not yet preserve.
